\documentclass[manuscript,screen]{acmart}

%%
%% \BibTeX command to typeset BibTeX logo in the docs
\AtBeginDocument{%
  \providecommand\BibTeX{{%
    \normalfont B\kern-0.5em{\scshape i\kern-0.25em b}\kern-0.8em\TeX}}}

%% Rights management information.  This information is sent to you
%% when you complete the rights form.  These commands have SAMPLE
%% values in them; it is your responsibility as an author to replace
%% the commands and values with those provided to you when you
%% complete the rights form.
\setcopyright{acmcopyright}
\copyrightyear{2026}
\acmYear{2026}
\acmDOI{XXXXXXX.XXXXXXX}

%% These commands are for a JOURNAL article.
\acmJournal{TOMS}
\acmVolume{1}
\acmNumber{1}
\acmArticle{1}
\acmMonth{2}

%%
%% Submission ID.
%% Use this when submitting an article to a sponsored event. You'll
%% receive a unique submission ID from the organizers
%% of the event, and this ID should be used as the parameter to this command.
%%\acmSubmissionID{123-A56-BU3}

%%
%% For managing citations, it is recommended to use bibliography
%% files in BibTeX format.
%%
%% You can then either use BibTeX with the ACM-Reference-Format style,
%% or BibLaTeX with the acmauthoryear style.
%%
\usepackage{listings}
\usepackage{xcolor}

\lstdefinelanguage{Julia}{
  morekeywords={abstract,break,case,catch,const,continue,do,else,elseif,end,export,false,for,function,immutable,import,importall,if,in,macro,module,quote,return,switch,true,try,type,typealias,using,while},
  sensitive=true,
  morecomment=[l]{\#},
  morecomment=[n]{\#=}{=\#},
  morestring=[b]",
  morestring=[b]',
  morestring=[b]"""
}

\lstset{
    language=Julia,
    basicstyle=\ttfamily\small,
    keywordstyle=\bfseries\color{blue},
    commentstyle=\itshape\color{green!40!black},
    stringstyle=\color{red},
    numbers=left,
    numberstyle=\tiny\color{gray},
    stepnumber=1,
    numbersep=5pt,
    backgroundcolor=\color{white},
    showspaces=false,
    showstringspaces=false,
    showtabs=false,
    frame=single,
    rulecolor=\color{black},
    tabsize=2,
    breaklines=true,
    breakatwhitespace=true,
    title=\lstname,
}

%%
%% end of the preamble, start of the body of the document source.
\begin{document}

%%
%% The "title" command has an optional parameter,
%% allowing the author to define a "short title" to be used in page headers.
\title{IntU.jl: A Julia Package for Symbolic Integration over Haar Measures}

%%
%% The "author" command and its associated commands are used to define
%% the authors and their affiliations.
%% Of note is the shared affiliation of the first two authors, and the
%% "authornote" and "authornotemark" commands
%% used to denote shared contribution to the research.
\author{\L ukasz Pawela}
\email{lpawela@iitis.pl}
\affiliation{%
  \institution{Institute of Theoretical and Applied Informatics, Polish Academy of Sciences}
  \streetaddress{Ba{\l}tycka 5}
  \city{Gliwice}
  \postcode{44-100}
  \country{Poland}
}

\author{Zbigniew Pucha\l a}
\email{z.puchala@iitis.pl}
\affiliation{%
  \institution{Institute of Theoretical and Applied Informatics, Polish Academy of Sciences}
  \streetaddress{Ba{\l}tycka 5}
  \city{Gliwice}
  \postcode{44-100}
  \country{Poland}
}

%%
%% By default, the full list of authors will be used in the page
%% headers. Often, this list is too long, and will overlap
%% with the page headers. This command allows the author to define
%% a more concise list of authors' names for this purpose.
\renewcommand{\shortauthors}{Pawela and Pucha\l a}

%%
%% The abstract is a short summary of the work to be presented in the
%% article.
\begin{abstract}
Symbolic integration over the Haar measure of compact groups is a computational
cornerstone in quantum information science and random matrix theory. We present
\texttt{IntU.jl}, a comprehensive Julia package for computing exact expectations
of polynomial functions over the unitary $U(d)$, orthogonal $O(d)$, and
symplectic $Sp(d)$ groups, as well as random pure quantum states. Unlike
existing tools that often rely on proprietary engines or lack support for
symbolic parameters, \texttt{IntU.jl} provides a fully open-source, performant
implementation of the Weingarten calculus. Key features include the native
handling of symbolic dimensions $d$ as rational functions, allowing for the
automatic derivation of asymptotic expansions, and a high-level symbolic trace
interface that reconstructs Weingarten graphs directly from index-free
expressions. We discuss the underlying algorithms—including the Murnaghan-Nakayama
rule and symplectic-orthogonal duality—and demonstrate that the package
efficiently handles high-degree moments, significantly outperforming manual derivations
and enabling new avenues of research in high-dimensional quantum systems.
\end{abstract}

%%
%% The code below is generated by the tool at http://dl.acm.org/ccs.cfm.
%% Please copy and paste the code instead of the example below.
%%
\begin{CCSXML}
<ccs2012>
   <concept>
       <concept_id>10011007.10011006.10011066.10011069</concept_id>
       <concept_desc>Software and its engineering~Software libraries and
repositories</concept_desc>
       <concept_significance>500</concept_significance>
       </concept>
   <concept>
       <concept_id>10002950.10003714.10003715</concept_id>
       <concept_desc>Mathematics of computing~Mathematical
software</concept_desc>
       <concept_significance>500</concept_significance>
       </concept>
 </ccs2012>
\end{CCSXML}

\ccsdesc[500]{Software and its engineering~Software libraries and repositories}
\ccsdesc[500]{Mathematics of computing~Mathematical software}

%%
%% Keywords. The author(s) should pick 3--5 keywords to be treated as
%% keywords in each piece of poetry together with their running enumeration.
\keywords{Symbolic integration, Haar measure, Weingarten calculus, Julia programming language, quantum information.}

%%
%% This command processes the author and affiliation and title
%% information and builds the first part of the formatted document.
\maketitle

\section{Introduction}
\label{sec:intro}
Integration over the Haar measure of compact groups is a cornerstone technique
in mathematical physics, with profound implications in random matrix theory
(RMT) \cite{mehta2004random, forrester2010log, livan2018introduction} and
quantum information theory (QIT) \cite{nielsen2010quantum, watrous2018theory}.
The ability to compute moments of random matrices enables the exact
characterization of ``typical'' properties of quantum systems, such as the
average entanglement entropy of subsystems—a concept famously pioneered by Page
\cite{page1993average}. Furthermore, these integrals are essential for studying
quantum chaos, scrambling, and the design of quantum circuits that approximate
randomness, known as unitary $t$-designs \cite{brandao2016exponential}.

Historically, the evaluation of these integrals was a formidable challenge.
While the defining properties of Haar measure were established by Hurwitz and
Weyl in the early 20th century, a systematic method for evaluating polynomial
integrals emerged later with the work of Weingarten
\cite{weingarten1978asymptotic}. He showed that asymptotic integrals could be
expressed as sums over permutations involving a weight function now called the
Weingarten function. This theory was rigorously detailed and extended to finite
$d$ by Collins \cite{collins2003moments}. The seminal work of Collins and Śniady
\cite{collins2006integration} unified the treatment for the unitary, orthogonal,
and symplectic groups, providing explicit combinatorial formulas for the
Weingarten functions in terms of characters of the symmetric group (or related
structures).

Despite the existence of these closed-form solutions, their practical evaluation
remains non-trivial due to the combinatorial complexity of the symmetric group
sums and the intricacies of representation theory. Several software packages
have been developed to automate these calculations. Notable examples include
\texttt{RTNI} \cite{fukuda2019rtni}, a Mathematica package that offers robust
functionality but relies on the proprietary Wolfram engine, and \texttt{Haarpy}
\cite{cardin2024haarpy}, a Python library that provides accessible symbolic
integration. Earlier work by Puchała and Miszczak \cite{puchala2017symbolic}
also explored symbolic integration algorithms.

However, existing tools often face limitations in extensibility, performance, or
the ability to handle symbolic dimensions $d$ natively as rational functions.
Furthermore, the translation from high-level, coordinate-free mathematical
expressions (e.g., traces of products) to the index-based tensor contractions
required by Weingarten calculus is often left to the user, leading to
error-prone manual preprocessing.

While \texttt{Haarpy} supports a broader range of groups—including circular
ensembles (COE, CSE) and permutation groups, which \texttt{IntU.jl} currently
does not, our package focuses on maximizing performance and symbolic flexibility
within the Unitary, Orthogonal, and Symplectic groups. A primary distinction
lies in \texttt{IntU.jl}'s deep integration with the Julia \texttt{Symbolics.jl}
ecosystem, treating dimensions $d$ as native symbolic rational functions rather
than requiring separate symbolic algebra backends or manual substitutions.

In this work, we present \texttt{IntU.jl}, a comprehensive Julia package
\cite{bezanson2017julia} designed to address these challenges. \texttt{IntU.jl}
offers a pure Julia implementation of the Weingarten calculus for the Unitary
$U(d)$, Orthogonal $O(d)$, and Symplectic $Sp(d)$ groups. Key features
distinguishing our approach include:
\begin{itemize}
    \item \textbf{Symbolic dimensions}: Fully symbolic treatment of the dimension
    $d$, enabling the derivation of exact asymptotic expansions and universal scaling laws.
    \item \textbf{Symbolic trace logic}: A high-level interface that allows users to
    integrate trace polynomials directly, with the software automatically handling
    the reduction to Weingarten graphs.
    \item \textbf{Performance}: Leveraging Julia's just-in-time compilation and
    efficient caching strategies to handle high-degree moments.
\end{itemize}

The remainder of this paper is organized as follows. Section
\ref{sec:background} reviews the mathematical background of Haar integration.
Section \ref{sec:features} describes the software architecture and core
features. Section \ref{sec:examples} provides usage examples, and Section
\ref{sec:implementation} details the implementation. We present benchmarks in
Section \ref{sec:performance} and conclude in Section \ref{sec:conclusion}.

\section{Background: Weingarten calculus}
\label{sec:background}
The core technique behind \texttt{IntU.jl} is the Weingarten calculus. For a
compact group $G$, the integral of a polynomial in matrix elements can be
expressed as a sum over a combinatorial set—permutations for the unitary case
and pair partitions for the orthogonal and symplectic cases.

\subsection{The unitary group $U(d)$}
For the unitary group $U(d)$, the integral of a polynomial of degree $k$ in both
$U$ and $\bar{U}$ is given by:
\begin{equation}
\int_{U(d)} U_{i_1 j_1} \dots U_{i_k j_k} \bar{U}_{i'_1 j'_1} \dots \bar{U}_{i'_k j'_k} dU = \sum_{\sigma, \tau \in S_k} \delta_\sigma(\vec{i}, \vec{i}') \delta_\tau(\vec{j}, \vec{j}') \mathrm{Wg}^U(\sigma\tau^{-1}, d),
\end{equation}
where $S_k$ is the symmetric group of degree $k$, and $\delta_\sigma(\vec{i},
\vec{i}')$ is a product of Kronecker deltas $\prod_m \delta_{i_m,
i'_{\sigma(m)}}$. The Weingarten function $\mathrm{Wg}^U(\sigma, d)$ depends
only on the cycle type of $\sigma$ and can be expressed via the irreducible
characters $\chi_\lambda$ of $S_k$:
\begin{equation}
\mathrm{Wg}^U(\sigma, d) = \frac{1}{(k!)^2} \sum_{\lambda \vdash k} \frac{f^\lambda \chi_\lambda(\sigma)}{s_{\lambda, d}(1)},
\end{equation}
where $f^\lambda$ is the dimension of the irreducible representation $\lambda$
and $s_{\lambda, d}(1)$ is the Schur polynomial evaluated at the identity,
representing the dimension of the corresponding $U(d)$ representation.

Crucially, $s_{\lambda, d}(1)$ is a polynomial in $d$, which implies that
$\mathrm{Wg}^U(\sigma, d)$ is a rational function of $d$. This property is
essential for symbolic integration, as it allows for exact results even when $d$
is treated as a variable. The function contains poles at integers $d$ such that
$|d| < k$, reflecting the breakdown of the expansion for small dimensions.

In the asymptotic limit $d \to \infty$, the Weingarten function behaves as:
\begin{equation}
\mathrm{Wg}^U(\sigma, d) \sim d^{-k - |\sigma|},
\end{equation}
where $|\sigma|$ denotes the minimum number of transpositions required to
generate $\sigma$, related to the number of cycles $c(\sigma)$ by $|\sigma| =
k - c(\sigma)$. For the identity permutation, the leading order term is
$\mathrm{Wg}^U(\mathrm{id}, d) \approx d^{-k}$, recovering the normalization
expected from independent Gaussian entries at the naive limit.

\subsection{Orthogonal and symplectic groups}
For the orthogonal group $O(d)$ and symplectic group $Sp(d)$, the integration
involves even moments of the matrix elements and is indexed by the set of pair
partitions $P_{2k}$ of $\{1, \dots, 2k\}$. For $O(d)$, the integral is:
\begin{equation}
\int_{O(d)} O_{i_1 j_1} \dots O_{i_{2k} j_{2k}} dO = \sum_{p, q \in P_{2k}} \delta_p(\vec{i}) \delta_q(\vec{j}) \mathrm{Wg}^O(p, q, d).
\end{equation}
Similarly, for the symplectic group $Sp(d)$ (where $d$ must be even), the
formula is:
\begin{equation}
\int_{Sp(d)} S_{i_1 j_1} \dots S_{i_{2k} j_{2k}} dS = \sum_{p, q \in P_{2k}} (-1)^{\ell(p) + \ell(q)} \delta_p(\vec{i}) \delta_q(\vec{j}) \mathrm{Wg}^{Sp}(p, q, d),
\end{equation}
where the signs depend on the contraction structure relative to the symplectic
form $J$. For detailed closure relations on these groups, see Matsuki
\cite{matsuki1990closure}.

\texttt{IntU.jl} leverages the fundamental duality between the orthogonal and
symplectic Weingarten functions:
\begin{equation}
\mathrm{Wg}^{Sp}(p, q, d) = (-1)^{k + \mathrm{loops}(p, q)} \mathrm{Wg}^O(p, q, -d),
\end{equation}
where $\mathrm{loops}(p, q)$ is the number of loops in the graph formed by the
union of the two pair partitions $p$ and $q$. This relation allows the package
to unify the codebase, computing symplectic integrals by reusing the optimized
orthogonal engine with an appropriately signed dimension parameter $d \to -d$.

Both $\mathrm{Wg}^O$ and $\mathrm{Wg}^{Sp}$ are rational functions of $d$, and
their poles dictate the dimensions for which the integration formulas are valid
(e.g., $d \ge k-1$ for $O(d)$).

\texttt{IntU.jl} efficiently computes these functions using both character-based
formulas and recursive relations, such as the Samuel-Weingarten recursion,
caching results to accelerate repeated integration tasks.

\section{Software features}
\label{sec:features}
\texttt{IntU.jl} is built on top of the \texttt{Symbolics.jl} ecosystem
\cite{symbolicsjl}, providing a seamless experience for Julia users.

\subsection{Groups and measures}
The package supports a wide range of integration measures, grouped as follows:
\begin{itemize}
    \item \textbf{Compact Groups}: Standard Haar measures on the unitary $U(d)$,
    orthogonal $O(d)$, and symplectic $Sp(d)$ groups. These are defined via
    \texttt{dU(U, d)}, \texttt{dO(O, d)}, and \texttt{dSp(S, d)}, where the
    first argument is a symbolic matrix variable (constructed via
    \texttt{Symbolics.jl} or the \texttt{@symbolic\_dimension} macro) and the
    second is the dimension, which can itself be a symbolic variable.
    \item \textbf{Quantum States}: The induced measure on pure states
    $|\psi\rangle$ drawn from the complex projective space $\mathbb{C}P^{d-1}$
    \cite{zyczkowski2001induced, bengtsson2017geometry}, accessed via
    \texttt{dPsi(psi, d)}. This corresponds to the distribution of the first
    column of a Haar-random unitary matrix.
    \item \textbf{Gaussian Ensembles}: Measures for the Gaussian unitary (GUE),
    orthogonal (GOE), and symplectic (GSE) ensembles (\texttt{dGUE},
    \texttt{dGOE}, \texttt{dGSE}). These rely on Wick's theorem for integration
    rather than Weingarten calculus.
    \item \textbf{Unitary Designs}: Unitary $t$-designs (\texttt{dDesign}),
    which mimic the first $t$ moments of the Haar measure, useful for studying
    pseudo-randomness in quantum circuits.
\end{itemize}

\subsection{Gaussian random matrix ensembles}
In addition to compact groups, \texttt{IntU.jl} supports integration over the
Gaussian ensembles (GUE, GOE, GSE). These are implemented using Wick's theorem
\cite{wick1950evaluation} (Isserlis' theorem) for moment contraction. For a
Gaussian Hermitian matrix $H$, the expectation values are determined by the
pairwise contractions:
\begin{itemize}
    \item \textbf{GUE}: $\langle H_{ij} \bar{H}_{kl} \rangle = \delta_{il} \delta_{jk}$, leading to $\langle \mathrm{tr}(H^2) \rangle = d^2$.
    \item \textbf{GOE}: $\langle H_{ij} H_{kl} \rangle = \delta_{ik} \delta_{jl} + \delta_{il} \delta_{jk}$ (for real symmetric $H$), leading to $\langle \mathrm{tr}(H^2) \rangle = d^2 + d$.
    \item \textbf{GSE}: For self-dual Hermitian matrices, the package utilizes
    the duality relation:
    \begin{equation}
    \langle \mathrm{tr}(H^k) \rangle_{\mathrm{GSE}}(d) = (-1)^{\frac{k}{2} + 1} \langle \mathrm{tr}(H^k) \rangle_{\mathrm{GOE}}(-d),
    \end{equation}
    which implies $\langle \mathrm{tr}(H^2) \rangle = d^2 - d$.
\end{itemize}
This approach allows \texttt{IntU.jl} to compute moments for all three Beta
ensembles ($\beta=1, 2, 4$) while maintaining full support for symbolic
dimensions $d$.

\subsection{Unitary designs}
Unitary $t$-designs \cite{dankert2009exact, gross2007evenly,
ambainis2007quantum} are ensembles of unitary matrices that mimic the Haar
measure up to the $t$-th moment. Specifically, a set of unitaries $\mathcal{D}
\subset U(d)$ is a $t$-design if:
\begin{equation}
\mathbb{E}_{U \in \mathcal{D}} [P(U, \bar{U})] = \int_{U(d)} P(U, \bar{U}) dU
\end{equation}
for all polynomials $P$ of degree $q \le t$ in the entries of $U$ and $\bar{U}$.

\texttt{IntU.jl} provides the \texttt{dDesign(U, d, t)} measure to represent
such ensembles. The implementation strictly enforces the moment-matching
condition:
\begin{enumerate}
    \item For integrals of degree $q \le t$, the package returns the exact
    Haar-averaged result using standard Weingarten calculus.
    \item For integrals of degree $q > t$, the package raises an explicit error,
    preventing incorrect assumptions about the design's higher-moment behavior.
\end{enumerate}
This feature allows researchers to verify whether specific quantum protocols or
randomized benchmarking schemes rely only on the $t$-moment properties of the
sampling distribution.

\subsection{Quantum information utilities}
The package includes a suite of helper functions tailored for quantum information
tasks, designed to work seamlessly with symbolic expressions:
\begin{itemize}
    \item \textbf{Purity}: The \texttt{purity($\rho$)} function computes
    $\mathrm{tr}(\rho^2)$. The companion function \texttt{average\_purity}
    automates the integration of this quantity over the specified measure,
    returning analytical averages for random states or channels.
    \item \textbf{Fidelity}: The \texttt{fidelity($\rho, \sigma$)} function
    calculates the fidelity between two states. For pure states or when one state
    is pure, this simplifies to the Hilbert-Schmidt inner product $\mathrm{tr}(\rho
    \sigma)$. \texttt{average\_fidelity} computes the expectation value of this
    metric.
    \item \textbf{Partial Trace}: The \texttt{partial\_trace(M, dims, system)}
    function computes the partial trace of a symbolic matrix $M$ over a specified
    subsystem. It handles composite dimensions $d = d_A d_B$ symbolically, ensuring
    that the resulting reduced density matrix retains the correct symbolic
    dependencies for further integration.
\end{itemize}
These utilities reduce the boilerplate code required for setting up common
calculations in quantum information theory, such as entanglement quantification
and randomized benchmarking.

\subsection{Symbolic dimensions and asymptotics}
A standout feature of \texttt{IntU.jl} is its deep support for symbolic dimensions $d$.
By leveraging the polynomial nature of group characters and Schur functions,
\texttt{IntU.jl} computes integrals as exact rational functions of $d$:
\begin{equation}
\int_{G(d)} P(U, \bar{U}) dU = \frac{N(d)}{D(d)},
\end{equation}
where $N(d)$ and $D(d)$ are polynomials. This capability is critical for studying:
\begin{itemize}
    \item \textbf{Transition points}: Identifying poles in $D(d)$ that correspond to
    dimensions where the Weingarten expansion breaks down (e.g., $d < k$).
    \item \textbf{Thermodynamic limits}: The \texttt{asymptotic(expr, measure, order)}
    function computes the series expansion of the result in powers of $1/d$.
    This automates the extraction of large-$d$ behavior, facilitating the study
    of concentration of measure and asymptotic freeness.
\end{itemize}
Even when the input dimension is numeric (e.g., $d=3$), the \texttt{asymptotic}
routine can introduce a dummy symbolic variable to perform the expansion, providing
theoretical insights alongside numerical results.

\subsection{Symbolic trace logic}
To simplify the integration of high-rank tensor networks and complex trace
expressions, \texttt{IntU.jl} introduces a \textbf{symbolic trace logic} system.
This system abstracts away explicit tensor indices, allowing users to define
computations in terms of coordinate-free matrix objects:
\begin{itemize}
    \item \textbf{Symbolic Matrix}: The \texttt{SymbolicMatrix} type represents
    an operator with specific properties (Unitary $U$, Adjoint $U^\dagger$, or
    Constant $M$).
    \item \textbf{Lazy Evaluation}: Operations like multiplication ($A * B$) and
    trace ($\mathrm{tr}$) generate \texttt{LazyTrace} objects. These objects maintain
    algebraic structures (products of traces of matrix strings) without expanding
    indices, e.g., representing $\mathrm{tr}(U A U^\dagger B)$ as a graph cycle rather
    than a sum over four indices.
    \item \textbf{Graph Reconstruction}: During integration, the package
    interprets these lazy structures as edges in a Weingarten graph. It
    automatically reconstructs the connectivity required for pair-partition
    matching, effectively converting the ``index-free'' input into the precise tensor
    contractions required by the Weingarten formula.
\end{itemize}
This abstraction significantly reduces the potential for index-mismatch errors
and renders the code nearly isomorphic to the mathematical pen-and-paper
formulation of the problem.

\section{Usage examples}
\label{sec:examples}
We illustrate the capabilities of \texttt{IntU.jl} with several examples, ranging
from basic component-wise integration to high-level symbolic trace manipulations
and asymptotic expansions. These examples demonstrate how the package's design
principles facilitate practical research tasks.

\subsection{Basic integration}
Consider the partial moment $\int_{U(d)} |U_{11}|^2 dU$. This integral represents
the average magnitude squared of a single entry of a Haar-random unitary matrix.
In \texttt{IntU.jl}, this can be computed using standard symbolic variables.
Note that we define a symbolic dimension $d$ and a symbolic unitary matrix $U$
of size $d \times d$:

\begin{lstlisting}[language=Julia]
using IntU, Symbolics
@variables d
@symbolic_dimension U[1:d, 1:d]
measure = dU(U)
expr = abs(U[1,1])^2
result = integrate(expr, measure)
# Output: 1/d
\end{lstlisting}

The result $1/d$ confirms the expected normalization: since the rows and columns
of a unitary matrix are unit vectors, the average squared magnitude of any entry
must be $1/d$.

\subsection{Symbolic Trace Logic}
For more complex expressions involving traces of random matrices, manual index
contraction is error-prone. The symbolic trace interface simplifies the
workflow by allowing users to define the calculation in coordinate-free notation.
Here, we compute $\int \mathrm{tr}(U A U^\dagger B) dU$, a standard integral
appearing in the study of quantum channels (specifically, the definition of the
depolarizing channel):

\begin{lstlisting}[language=Julia]
# Define symbolic matrices
U_var = SymbolicMatrix(:U)
A = SymbolicMatrix(:A); B = SymbolicMatrix(:B)
measure = dU(U_var, d)

# Compute integral of tr(U A U' B)
# tr_lazy constructs a lazy trace object
expr = tr_lazy(U_var * A * U_var' * B)
result = integrate(expr, measure)
# Output: tr(A)tr(B)/d
\end{lstlisting}

This result neatly recovers the identity $\mathbb{E}[U A U^\dagger] =
\mathrm{tr}(A) \mathbb{I} / d$.

\subsection{Asymptotic expansions}
Analyzing behavior in the large-$d$ limit is crucial for understanding
thermodynamic properties and concentration of measure. \texttt{IntU.jl}
facilitates this analysis by expanding exact rational results into power series.
Consider the fourth moment of a matrix entry, $|U_{11}|^4$:

\begin{lstlisting}[language=Julia]
expr = abs(U[1,1])^4
# Exact result computed internally: 2/(d*(d+1))
asymp_res = asymptotic(expr, measure, 4)
# Output: 2/d^2 - 2/d^3 + 2/d^4
\end{lstlisting}

\section{Implementation details}
\label{sec:implementation}
\texttt{IntU.jl} is written in pure Julia \cite{bezanson2017julia}, utilizing
the language's multiple dispatch and metaprogramming features to build a highly
modular and extensible system. The core architecture consists of three main
components: the integration pipeline, the Weingarten engine, and the symbolic
trace logic.

\subsection{Integration pipeline}
The integration process follows a systematic pipeline designed to normalize and
dispatch symbolic expressions efficiently:
\begin{enumerate}
    \item \textbf{Normalization}: The input expression is first normalized using a ruleset
    based on \texttt{SymbolicUtils.jl}. This step expands complex conjugates and rewrites
    composite functions—such as $\mathrm{abs}(z)^2 \to z \bar{z}$, $\mathrm{real}(z) \to \frac{1}{2}(z +
    \bar{z})$, and $\mathrm{imag}(z) \to \frac{1}{2i}(z - \bar{z})$—into polynomial forms
    amenable to integration.
    \item \textbf{Substitution}: Symbolic variables representing matrix elements
    $U_{ij}$ are replaced by an internal atomic representation that facilitates
    pattern matching against the integration measure.
    \item \textbf{Expansion}: The expression is algebraically expanded into a sum of monomials.
    This relies on the efficient sparse polynomial handling of
    \texttt{Symbolics.jl}.
    \item \textbf{Dispatch}: Each monomial is analyzed to extract the indices of
    the random matrices. The integration is then dispatched to the specific
    algorithm appropriate for the group ($U, O, Sp$) or measure (Gaussian), relying on
    Julia's multiple dispatch to handle different measure types without branching logic.
\end{enumerate}

\subsection{Weingarten engine}
The calculation of Weingarten functions is computationally intensive. To achieve
high performance, \texttt{IntU.jl} implements specialized algorithms for each group:
\begin{itemize}
    \item \textbf{Unitary Group}: Characters of the symmetric group are computed
    using the Murnaghan-Nakayama rule, which is significantly faster than determinantal
    formulas for sparse partitions. Dimensions of irreducible representations $s_{\lambda, d}(1)$
    are computed via the Hook-Content formula:
    \begin{equation}
    f^\lambda_d = \prod_{(i,j) \in \lambda} \frac{d + c_{i,j}}{h_{i,j}}
    \end{equation}
    where $c_{i,j} = j-i$ is the content and $h_{i,j}$ is the hook length. This formulation
    allows the package to handle symbolic dimensions natively, producing exact rational functions
    of $d$.
    \item \textbf{Orthogonal/Symplectic Groups}: For these groups, the package generates pair partitions
    recursively, canonicalizing them (typically as sorted pairs) to enable $\mathcal{O}(1)$ lookup.
    It constructs the Gram matrix $G_{\pi, \sigma} = d^{\ell(\pi, \sigma)}$, where $\ell(\pi, \sigma)$
    is the number of loops in the graph formed by pair partitions $\pi$ and $\sigma$. The formulation
    inverts this matrix analytically. Crucially, symplectic Weingarten functions are evaluated
    by reusing the orthogonal engine via duality:
    \begin{equation}
    \mathrm{Wg}^{Sp}(p, q, d) = (-1)^{k + \mathrm{loops}(p, q)} \mathrm{Wg}^O(p, q, -d),
    \end{equation}
    avoiding the need for a separate symplectic solver.
    \item \textbf{Memoization}: All computed characters, dimensions, and
    Weingarten values are cached using the \texttt{Memoize.jl} package. Memoization occurs
    at the granular level of individual character and dimension calls, ensuring maximizing
    reuse across different partition structures and integration orders.
\end{itemize}

\subsection{Symbolic Trace Logic}
To handle high-level expressions, \texttt{IntU.jl} introduces specialized
abstract types: \texttt{LazyTrace} and \texttt{LazySum}. These structures maintain
the algebraic form of trace products (e.g., $\mathrm{tr}(A B) \mathrm{tr}(C)$)
without strictly evaluating them into scalar components.

The integration strategy for these structures is twofold:
\begin{enumerate}
    \item \textbf{Library Lookup}: Common patterns, such as the single-channel
    overlap $\mathrm{tr}(U A U^\dagger B)$, are detected via pattern matching
    mechanisms (implemented in `check\_library`). These matches trigger optimized
    fast-paths that return pre-computed results directly (e.g.,
    $\frac{1}{d}\mathrm{tr}(A)\mathrm{tr}(B)$) without invoking the full Weingarten engine.
    This effectively reduces the complexity for these common cases from $O((k!)^2)$ to $O(1)$.
    \item \textbf{Graph Reconstruction}: For generic cases, the package
    interprets the trace cycles as edges in a graph. By identifying the relative
    positions of $U$ and $U^\dagger$ within the lazy structure, the integrator
    assigns effective indices to the intervening constant scaling matrices.
    This effectively converts the ``index-free'' input into the precise tensor
    contraction graph required by the Weingarten formula calculation, thereby
    avoiding the combinatorial explosion associated with eager index expansion.
\end{enumerate}

\section{Performance and benchmarks}
\label{sec:performance}
We evaluated the performance of \texttt{IntU.jl} using a suite of stress tests
involving high-degree polynomial integrals over the Unitary, Orthogonal, and
Symplectic groups. The benchmarks were run on a standard workstation (Julia 1.12.3).
Specifically, we evaluated the computation time for the moments of single matrix entries,
such as $\int |U_{11}|^{2k} dU$, for high degrees $2k \geq 8$. These integrals represent
a standard stress test for Weingarten implementations due to their maximal combinatorial
complexity.

All benchmarks were performed using the \texttt{BenchmarkTools.jl} package.
Reported execution times represent the \textbf{median} of $N=30$ runs.
Crucially, these results \textbf{exclude} the initial JIT compilation and
warm-up time, reflecting the true steady-state performance of the library
observed during production use. The code for these specific stress tests is
located in \texttt{benchmarks/09\_stress\_test\_2.jl}, which can be executed via
the reproduction script described in Section \ref{sec:availability}. Additional
benchmarks for lower-degree moments and application examples are available in
\texttt{benchmarks/08\_stess\_test\_1.jl}.



\subsection{Theoretical complexity}
The computational complexity of the Weingarten integration is dominated by the
summation over the symmetric group $S_k$ (for the unitary case) or the set of
pair partitions $M_{2k}$ (for orthogonal/symplectic cases). The size of these
sets grows factorially: $|S_k| = k!$ and $|M_{2k}| = (2k-1)!!$. Consequently,
the runtime scales exponentially with the degree of the polynomial.
Empirically, we find that exact symbolic integration is feasible for degrees up
to $2k \approx 10-12$ on standard hardware. Beyond this regime, the number of
terms (e.g., $10! \approx 3.6 \times 10^6$) becomes prohibitive for real-time
symbolic manipulation, although numerical evaluation remains tractable for
slightly higher degrees.

\subsection{Benchmark results}
Table \ref{tab:benchmarks} summarizes the execution times for selected integrals.
Key observations include:
\begin{enumerate}
    \item \textbf{Symbolic overhead}: Computing integrals with symbolic dimension
    $d$ is computationally more demanding than fixed $d$, particularly for the
    unitary group. For example, computing $\int |U_{11}|^{10} dU$ symbolically
    took approximately 6.9 seconds, whereas the same integral for fixed $d=50$
    completed in 0.18 seconds. This overhead arises from the manipulation of
    large rational polynomials in $d$.
    \item \textbf{Low-degree performance}: For common use cases involving
    lower-degree polynomials (e.g., $2k=6$), the package is extremely fast. The
    symbolic integration of $|U_{11}|^6$ completes in approximately 8.9 ms, and
    for the orthogonal group, computing $O_{11}^6$ takes only 0.3 ms for $d=10$.
    \item \textbf{Group differences}: Orthogonal and Symplectic integrals are
    generally faster to compute than Unitary ones of comparable degree. This is
    typically due to the smaller size of the summation formulations (pair partitions vs permutations)
    at moderate degrees, although the number of pair partitions grows faster asymptotically.
    \item \textbf{Correctness}: All benchmarks were verified against known exact
    values or asymptotic results, confirming the numerical stability of the
    implementation even for large dimensions (e.g., $d=50$).
    \item \textbf{Application}: The purity calculation for a bipartite state
    ($d=6$), which involves contracting a 4-th degree polynomial over the
    unitary group, completes in approximately 81.1 ms, demonstrating the tool's
    applicability to quantum information tasks.
\end{enumerate}

\begin{table}[ht]
\centering
\caption{Benchmark results for symbolic vs. numeric integration. Times are median values over 30 samples (except for long-running symbolic cases).}
\label{tab:benchmarks}
\begin{tabular}{llcc}
\hline
Group & Integrand & Dimension $d$ & Time (ms) \\
\hline
Unitary & $|U_{11}|^6$ & Symbolic & 8.9 \\
Unitary & $|U_{11}|^8$ & Symbolic & 156.3 \\
Unitary & $|U_{11}|^{10}$ & Symbolic & 6912.9 \\
Unitary & $|U_{11}|^{10}$ & $d=10$ & 193.5 \\
Unitary & $|U_{11}|^{10}$ & $d=50$ & 183.6 \\
\hline
Orthogonal & $O_{11}^6$ & $d=10$ & 0.3 \\
Orthogonal & $O_{11}^8$ & $d=20$ & 2.9 \\
Orthogonal & $O_{11}^{10}$ & $d=20$ & 230.6 \\
Orthogonal & $O_{11}^{10}$ & $d=50$ & 244.4 \\
\hline
Symplectic & $|S_{11}|^8$ & $d=10$ & 2.1 \\
Symplectic & $|S_{11}|^{10}$ & $d=10$ & 4.9 \\
Symplectic & $|S_{11}|^{10}$ & $d=20$ & 9.9 \\
\hline
Application & Purity & $d=6$ & 81.1 \\
\hline
\end{tabular}
\end{table}

The results demonstrate that \texttt{IntU.jl} remains efficient for practically
relevant problem sizes. The rapid execution of the purity application ($\approx$
81 ms) highlights the library's utility for routine quantum information tasks,
while the sub-millisecond performance for low-degree orthogonal integrals
facilitates large-scale statistical sampling. Simultaneously, the symbolic
engine proves capable of handling high-degree moments that are intractable by
hand. The caching mechanism ensures that subsequent calls with the same
parameters are virtually instantaneous.

\section{Conclusion}
\label{sec:conclusion}
We have presented \texttt{IntU.jl}, a robust and extensible Julia package for
performing symbolic integration over the Haar measure of the Unitary, Orthogonal,
and Symplectic groups. By bridging the gap between the abstract theory of
Weingarten calculus and practical computation, \texttt{IntU.jl} empowers
researchers to rigorously derive properties of random quantum channels, higher-order
entanglement statistics, and scrambling measures without resorting to brittle or
approximate numerical sampling.

The package's core strengths lie in its unified treatment of compact groups,
its native handling of symbolic dimensions via rational functions, and its
novel symbolic trace logic. The latter effectively decouples the user's high-level
mathematical intent from the low-level index contractions required by the integration
engine, significantly reducing the cognitive load and potential for error. As
demonstrated by the benchmarks, the optimized Julia implementation ensures that
this symbolic abstraction does not come at the cost of performance, enabling the
evaluation of high-degree moments that were previously intractable.

Future development of \texttt{IntU.jl} will focus on expanding support to other
compact Lie groups, integrating more seamlessly with tensor network libraries
such as \texttt{ITensors.jl}, and further optimizing the asymptotic expansion
engine for extremely large-scale problems. We believe \texttt{IntU.jl} will become
an essential utility in the toolkit of theoretical physicists and quantum
information scientists.

\section{Software availability and reproducibility}
\label{sec:availability}
The source code for \texttt{IntU.jl} is available on GitHub under the Apache License 2.0.
\begin{itemize}
    \item \textbf{Repository URL}: \url{https://github.com/iitis/IntU}
    \item \textbf{License}: Apache License 2.0
    \item \textbf{Version}: The results in this paper were generated using
    version v0.3.0.
    \item \textbf{Benchmarks}: To reproduce the benchmarks presented in Section
    \ref{sec:performance}, execute the provided script from the repository root:
    \begin{lstlisting}[language=bash]
    bash benchmarks/runbenchmarks.sh
    \end{lstlisting}
\end{itemize}


\begin{acks}
LP and ZP acknowledge financial support from the National Science Center,
grant number\ldots
\end{acks}

\bibliographystyle{ACM-Reference-Format}
\bibliography{references}

\end{document}
\endinput